%==============================================================================
\chapter{Schedule}\label{cap:implementacao}
%==============================================================================

This chapter presents the schedule for the continuation of this work in the
Final Undergraduate Project~II (TCC~II). The present stage (TCC~I) established
the research questions, the experimental methodology, and a first round of
results covering five LLMs across six PAMPAR programs, using a single
zero-shot prompt with the sequential source code as input. The activities
planned for TCC~II extend this foundation along the directions identified in
Chapter~\ref{cap:conclusao}, deepening the evaluation and consolidating the
findings into a complete study.

Two of the planned directions deserve particular emphasis, as they broaden the
scope of the central research question from \textit{how well current LLMs
parallelize code} to \textit{how the way the task is requested affects the
result}.

The first direction concerns \textbf{prompt engineering}. The results obtained
in TCC~I used a single zero-shot prompt, in which the task is described without
worked examples. The literature on prompting describes several techniques that
can substantially change the quality of generated code. In few-shot prompting,
a small number of input--output examples are provided before the target task,
allowing the model to infer the expected pattern from
demonstrations~\citeonline{brown2020fewshot}. In chain-of-thought prompting, the
model is guided to reason explicitly through intermediate steps before producing
the final answer, which has been shown to improve performance on tasks that
require multi-step reasoning~\citeonline{wei2022cot}; a simple zero-shot variant
of this technique prepends an instruction such as ``let us reason step by
step''~\citeonline{kojima2022zeroshot}. These techniques are directly relevant
to the failures observed in TCC~I, several of which stemmed from a model
selecting an inappropriate parallelization strategy rather than from an inability
to write valid code. Evaluating whether guided prompting helps the models reason
about data dependencies before generating code is therefore a natural extension.

The second direction concerns \textbf{varying the input given to the model}. In
TCC~I, each model received the complete sequential source code and was asked
only to parallelize it, which isolates the parallelization capability from the
implementation capability. An alternative is to provide only the name or a
high-level description of the algorithm (for example, ``write a parallel OpenMP
implementation of matrix multiplication'') and require the model to implement
\emph{and} parallelize it. Comparing the two input conditions---with and without
the sequential code as a worked model---would reveal whether the models depend on
the provided implementation or are able to produce efficient parallel code from a
specification alone. This comparison also connects to a finding from TCC~I, in
which one model improved the algorithm itself rather than merely parallelizing
it; providing only a description might lead more models to such reimplementations.
A methodological challenge of this direction is correctness verification: when the
model implements the algorithm from scratch, ensuring that its output is
comparable to the sequential oracle requires additional care in defining inputs
and expected results.

The planned activities are:

\begin{alineas}
\item \textbf{A1 --- Literature review extension:} expand the review of related
work on LLM-based parallel code generation and prompt engineering, incorporating
newly published studies and positioning the present work within them.
\item \textbf{A2 --- Prompt engineering experiments:} evaluate different
prompting strategies (zero-shot, few-shot, and chain-of-thought) applied
uniformly across all models and programs, in order to measure how the prompting
approach affects the correctness and efficiency of the generated code.
\item \textbf{A3 --- Input variation experiments:} compare two input conditions,
one providing the complete sequential code and one providing only the name or a
high-level description of the algorithm, to assess the models' dependence on a
worked implementation.
\item \textbf{A4 --- Benchmark extension:} extend the evaluation to additional
PAMPAR programs, particularly Jacobi Method (JA), Game of Life (GL), and
Gram-Schmidt (GS).
\item \textbf{A5 --- Consolidated analysis:} integrate all results, refine the
failure-mode characterization, and strengthen the comparative analysis across
models, prompting strategies, and input conditions.
\item \textbf{A6 --- Writing and revision:} write the additional chapters,
revise the full text, and prepare the final document.
\item \textbf{A7 --- Final presentation:} prepare and deliver the TCC~II
defense.
\end{alineas}

Table~\ref{tab:schedule} presents the distribution of these activities across
the five months planned for TCC~II. Each filled cell ($\bullet$) indicates that
the corresponding activity is scheduled for that month.

\begin{table}[htbp]
\centering
\caption{Planned schedule of activities for TCC~II.}
\label{tab:schedule}
\begin{tabular}{lccccc}
\toprule
\textbf{Activity} & \textbf{M1} & \textbf{M2} & \textbf{M3} & \textbf{M4} & \textbf{M5} \\
\midrule
A1 --- Literature review extension   & $\bullet$ & $\bullet$ &           &           &           \\
A2 --- Prompt engineering            & $\bullet$ & $\bullet$ & $\bullet$ &           &           \\
A3 --- Input variation               &           & $\bullet$ & $\bullet$ & $\bullet$ &           \\
A4 --- Benchmark extension           &           &           & $\bullet$ & $\bullet$ &           \\
A5 --- Consolidated analysis         &           &           &           & $\bullet$ & $\bullet$ \\
A6 --- Writing and revision          &           &           &           & $\bullet$ & $\bullet$ \\
A7 --- Final presentation            &           &           &           &           & $\bullet$ \\
\bottomrule
\end{tabular}
\fonte{The author.}
\end{table}